\sscp{Size of sample for establishing patterns}{02-04-01}
We have already noted that the Austronesian languages in our
target region are severely under documented,
with very few published grammars or dictionaries,
a few short grammar
and phonology sketches, and a few additional papers on other
languages about limited topics (e.g. negation, subjecthood, possession).
While for several languages we have access to more recent,
more reliable, and more complete data,
much of the earlier comparative discussion is based on word lists,
many of which are very limited (200 items or less),
and of varying quality and reliability.
The sample is often simply too small to make definitive statements
about regular sound correspondences or to have confidence in
some of the glosses.

\citet[8]{Klamer2019} quantifies it this way:

\begin{quote}
	“The data sets on which earlier classifications of languages
	in this region have been based are very limited and narrow.
	For example, of the 208 languages that Glottolog
	(Hammarström, Forkel, {\&} Haspelmath, n.d.) [Glottolog 3.0] lists as spoken
	in eastern Indonesia, less than 50{\%} (102) are represented in the major reference
	work on Austronesian reconstructions,
	the Austronesian Comparative Dictionary
	(ACD; Blust {\&} Trussel, n.d. [\citet{BlustTrussel2020}]).
	Of those, about 50{\%} (52) features in the database with less
	than 10 words, and 25 of them with only one (!) word.
	[{\ldots}] Clearly, affiliations based on so little data collected
	from such a vast insular region comprising thousands of islands
	leave a lot of room for uncertainties.” \citep[8]{Klamer2019}
\end{quote}

Ten words in a language,
or even a hundred,
is not necessarily going to reveal,
for instance, that PMP *R > \it{l} word initially,
but *R > \it{r} between vowels,
or that PMP *j > \it{l} in some words or environments,
but *j > {\0} in others in the same language.
With such a small data sample we have little basis by which
to gauge the accuracy of the transcription (e.g.
are glottal stops indicated? double vowels? schwa? is orthographic
‹b› fricative [β], implosive [ɓ]?
Is ‹f› bilabial [ɸ]?) or the reliability
and completeness of glosses (e.g.
Is the word for ‘sun’ also the word for ‘day’? Is the word for
‘boat’ generic, or refers only to ‘outrigger canoe’? Does the word given for
‘bird’ also include other winged creatures such as bats?).
Is the word for ‘grandparent’ a reciprocal term that also includes
‘grandchild’? Is the word glossed ‘brother’,
from the perspective of a same sex sibling,
opposite sex sibling, or either?

\ili{As} seen in the following chapters,
some comparative linguists are willing to draw conclusions
and make subgrouping arguments on the basis of historical medial
*NC consonant clusters, such as *mp, *mb,
*nt, *nd, *ns, *nz,
*ŋd, *md, *ŋk, *ŋg.
But the reality in the data available to us in our target region
shows a variety of problems that should encourage caution.

\begin{enumerate}[label=(\alph*)]
	\item	The number of proto forms for which we have data in which
				such sequences occur is quite low,
				often an occurrence of only one or none.
				This is true even for fairly well documented languages.
	\item	The number of words which reflect such medial clusters
				found in enough daughter languages to see a comparative pattern
				is also quite low to non-existent.
	\item	Some languages reflect the presence of historical clusters
				and others do not (e.g. *upu vs.
				*umpu ‘grandparent/child’, *ŋisi vs. *ŋinsi ‘tooth’).
				To illustrate a number of issues under discussion:
		\begin{itemize}
			\item In \citet{BlustTrussel2020}, \ili{Buru} \it{upu}
						is assigned to PMP *umpu ‘grandparent/ grandchild (recipr.); ancestor’.
			\item But \it{upu} is the form for the \ili{Lisela} language.
						The \ili{Buru} form is actually \it{opo}, and \it{opo=n}.
			\item \ili{Buru}
						and \ili{Lisela} sound correspondences are both *p = \it{p}, *mp > \it{b}.
						These are very consistent (except in a few loan words \ili{Lisela}
						has borrowed from \ili{Kayeli} which reflect *p > \it{h}).
			\item The expected form from *umpu for \ili{Buru} would be an unattested *[obo].
						So we can say with a high degree of confidence that \ili{Buru} \it{opo}
						is from a parent form of *upu, not *umpu.
						(Some nearby languages do reflect *umpu,
						but \ili{Buru} and many other languages in the same region reflect
						*upu; still others are not clear.)
		\end{itemize}
	\item	A single historical cluster may reflect three or four different
				patterns in the languages (see \citealp{Grimes1991a}).
				For example, *ma-diŋdiŋ ‘cold’ may be retained in different languages
				as \it{didi, diŋin, rini, marindin, marihi,
				ridi, riki}, or \it{lirin,} reflecting several different responses
				to the medial cluster.
	\item	A single language may reflect more than one response.
				For example, \ili{Kei} Kecil has both \it{ŋa-ridin}
				‘cold’ and \it{na-b-rinin} ‘chills’.
	\item	A single language may reflect an *NC cluster in one word,
				and a medial *C without the nasal in another word.
				For example PMP *upu yields the expected \ili{Buru} \it{opo,
				opo=n} ‘1) grandparent-grandchild (reciprocal term); 2) master, lord’.
				Whereas the related PMP *ta umpu ‘grandparent/grandchild (referential
				term used reciprocally)’ yields the expected \ili{Buru} form \it{tobo=n}
				‘person responsible, master, lord, one with custodial rights, owner, sir, host’.
				So care needs to be taken not to assume all reconstructed medial
				*NC clusters behave the same way in the same language for all words.
				This issue is troublesome for the whole idea of regular sound
				correspondences, unless one treats the medial *NC cluster as optional for many
				words in many languages,
				and checks each item separately.
	\item	\ili{As} above,
				two words in a language may reflect two different patterns for
				the same historical cluster.
				For example, the reflexes for *ns in *ŋinsi (local gloss: ‘tooth’)
				show a different pattern than for *mans(aə)r (local gloss: cuscus).
				There may be no other words found in that language reflecting
				*ns in a proto form.
				An occurrence of one is not a pattern
				and does not qualify as a regular sound correspondence.	
\end{enumerate}

We tend towards caution on making subgrouping claims based on
rare or elusive historical *CC clusters -- unless there is actually
sufficient data to establish a pattern
and substantiate a claim.
We prefer to see \it{a minimum} of two to three examples showing
the same pattern before allowing that it might reflect a regular sound correspondence.
We remain sceptical about claims of sound correspondences for
rare or elusive data reflecting *r,
*nj, *nz, *ns, *ŋk,
*ŋg, for example, since we ourselves can find very few unambiguous reflexes of these.
Even more rarely do we find enough reliable examples in related
languages to base claims about subgrouping on them.

What a small sample of ten words or so sometimes may allow,
is a glimpse of how it is similar to,
or differs from other languages in the same local region or node of classification.
The small sample may (or may not),
be enough to reflect a diagnostic feature for that node.

To take one example,
\citet[129]{Collins1986} classified \ili{Seti},
spoken in central-eastern Seram,
with \ili{Masiwang} and \ili{Bobot} partly on the basis of shared *ŋ > \it{k},
though he noted: “In \ili{Seti} the reconstruction is a problem because
so few cognates containing *ŋ occur.
Two display \it{n} while one shows \it{k}.” However,
examination of more data (to which Collins did not have access)
shows that the one word with *ŋ > \it{k} (*naŋuy > \it{nak} ‘swim’)
truly is an exception with ten other examples of *ŋ > \it{n}.
(See \srf{10-03-04}, particularly example
\ref{ex:10-59} for more discussion of \ili{Seti}).
\citet{Collins1986} classification was reasonable based on the
limited data available to him,
but it is illustrative of the problems
faced with many languages across the region.


